% Drawing 78/03-1 -- folder 78, batch 03, archive page 57 (PDF page 58).
% Transcription: Une ligne brisée~: le segment $a_0$ de $s_0$ à $s_1$, puis le segment $a_1$ de $s_1$ à $s_2$, puis l'amorce d'un troisième segment, vertical, au-dessus de $s_2$.
% Kind: broken line. Method 2 (schematic redraw from the facsimile). Status: ok.
% Notes: The segment a_0 from s_0 to s_1, the segment a_1 from s_1 to s_2, and the start of a vertical segment above s_2.
% Coordinates: 1 unit = 1 mm = 1/100 of the crop width (bbox in page fractions: [0.32, 0.51, 0.5, 0.6]).
% Redrawn by Claude (claude-opus-5-5), issue #73, 2026-11-07; not checked by a human; not compiled (no TeX in the session).
% Rights: reproduces part of the work -- behind RENDER_TRANSCRIPTS, never published by itself.
\documentclass[tikz,border=3pt]{standalone}
\usepackage{amsmath,amssymb}
\usetikzlibrary{patterns,arrows.meta}
\tikzset{pencil/.style={black!45}, struck/.style={black!35}, mark/.style={circle,draw,inner sep=1.3pt}, pt/.style={circle,fill,inner sep=0.9pt}, opt/.style={circle,draw,fill=white,inner sep=1.1pt}}
\begin{document}
\begin{tikzpicture}[x=1mm,y=1mm,line cap=round,line join=round,font=\scriptsize]
\draw (17.2,-53.8) -- (53.9,-50.5) -- (83.9,-28.8) -- (85.6,-5.5);
\node[pt] at (53.9,-50.5) {};
\node[pt] at (83.9,-28.8) {};
\node[inner sep=1pt] at (12.2,-60.5) {\tiny $s_0$};
\node[inner sep=1pt] at (58.9,-57.1) {\tiny $s_1$};
\node[inner sep=1pt] at (92.2,-33.8) {\tiny $s_2$};
\node[inner sep=1pt] at (33.9,-45.5) {\tiny $a_0$};
\node[inner sep=1pt] at (60.6,-35.5) {\tiny $a_1$};
\end{tikzpicture}
\end{document}
